\chapter{Proposal Feedback and Iteration Plan}

\section{Proposal Feedback Summary}\label{sec:feedback}
The proposal defence result was sent by the FYP Committee on 14 September 2026. The decision for
F26-155-GeoCrop was \textbf{Approved with Minor Modifications}. The panel's comments are summarised
below; each is numbered P1--P7 and tracked in the action register.
\begin{enumerate}[label=\textbf{P\arabic*.}, leftmargin=1.1cm]
  \item \textbf{Clarify the intended goals.} The proposal claimed crop mapping, GIS integration and
        real-time estimation without stating clearly what each meant.
  \item \textbf{Improve the evaluation metrics.} The proposed criteria did not measure the actual
        challenges of the project.
  \item \textbf{Measure the real challenges: computational performance, a lightweight system and
        usability.} The panel stated that the evaluation criteria were wrong, and that usability,
        performance and integration are the main concerns.
  \item \textbf{Analyse feasibility.} The complexity and resource requirements of the complete
        workflow had not been analysed beforehand.
  \item \textbf{Do not run everything in a single flow.} The computational and space requirements are
        high enough that running the entire process in a single flow may not be practical.
  \item \textbf{Integration of huge datasets.} Large datasets are not easy to load and use together.
  \item \textbf{Work on presentation skills.}
\end{enumerate}
The panel also asked the team to incorporate all changes it suggested. These comments set the agenda
for Iteration 1. The architecture track answered P3--P6 by redesigning the platform around
performance, light weight and integration, and measuring each one. The model track answered P2--P4
by changing the evaluation protocol and selecting the model on cost as well as accuracy.

The team's own review after the defence found further issues: the prototype could not run models on
uploaded images, the first evaluation used an optimistic random pixel split, models would keep
changing during the project, and the iteration plan in the proposal slides no longer matched the
order in which the work had to be done. These items are recorded separately in the register, with
their source.

\section{Proposal Action Register}
\begin{table}[htbp]
\centering\small
\caption{Proposal action register (summary; full register with evidence in Appendix~\ref{app:register})}
\label{tab:register}
\begin{tabularx}{\textwidth}{L{0.8cm} L{3.6cm} Y L{1.3cm} L{1.7cm}}
\toprule
\textbf{No.} & \textbf{Feedback / required change} & \textbf{Action taken} & \textbf{Status} & \textbf{Evidence}\\
\midrule
P1 & Clarify the intended goals. & Objectives restated with measurable targets; ``real-time'' defined as an uploaded image on the map in about a minute; scope split by FYP phase. & Done & Sec.~\ref{sec:goals}\\
P2, P3 & Evaluation metrics must measure performance, light weight, usability and integration. & Evaluation criteria rewritten around these concerns, each with a metric and target. Performance, cost and integration measured; usability checked against requirements, user study planned. & Partly done & Table~\ref{tab:evalcriteria}\\
P4 & Analyse feasibility and resource requirements of the whole workflow. & Time, memory and storage measured for every stage; the baseline crop-type model (3.1~GB RAM) found to dominate, so it is replaced by LightGBM (549~MB). & Done & Table~\ref{tab:resources}\\
P5 & A single-flow process is not practical. & Workflow split into independent queued stages; results kept as compressed rasters; maps served tile by tile on demand. & Done & Sec.~\ref{sec:feasibility}\\
P6 & Huge datasets are hard to load and use together. & Only visible tiles are sent: 337~KB per screen instead of 43.3~MB; layers share one catalogue. & Done & Table~\ref{tab:benchmark}\\
P7 & Presentation skills. & Fixed demonstration script on the running system; each member presents their own track; rehearsals before evaluations. & In progress & Sec.~\ref{sec:demo}\\
T1 & Team review: uploaded images were never analysed. & Upload, validation, preprocessing and model runs as background jobs; 57~s upload-to-map. & Done & TC-11, Fig.~\ref{fig:seq}\\
T2 & Team review: random pixel split overstates accuracy. & Spatial block split with buffer; choices on validation; intervals over test blocks. & Done & Sec.~\ref{sec:protocol}\\
T3 & Team review: models will keep changing. & Manifest-based model registry; add-a-model test. & Done & TC-03\\
T4 & Team review: iteration plan in the slides outdated. & Two-track iteration plan (Sections~\ref{sec:it1plan}, \ref{sec:it2plan}, Chapter~\ref{ch:remaining}). & Done & This chapter\\
\bottomrule
\end{tabularx}
\end{table}

\section{Iteration Strategy}
Development is iterative and time-boxed. Each iteration ends with a working, demonstrable increment,
a short evaluation, and an updated plan. Because GeoCrop has two kinds of risk, work proceeds on two
parallel tracks that meet at the end of each iteration:
\begin{itemize}
  \item \textbf{Architecture track} (platform): data flow, storage, processing, map delivery and
        interface. Risk: performance and integration.
  \item \textbf{Model track} (research): data preparation, model comparison and evaluation. Risk:
        accuracy, generalisation and cost.
\end{itemize}
The two tracks share one contract: a model is a plug-in that reads an aligned image and writes a
georeferenced raster, so research results can move into the platform without redesign.
Decisions are made on measured evidence (benchmarks for the platform, validation scores for models)
and recorded in this report and the repository.

\begin{table}[htbp]
\centering\small
\caption{Updated iteration plan (replaces the phase plan presented in the proposal slides)}
\label{tab:iterations}
\begin{tabularx}{\textwidth}{L{2.4cm} L{2.6cm} Y}
\toprule
\textbf{Iteration} & \textbf{Period} & \textbf{Focus}\\
\midrule
Iteration 1 & Sep -- mid-Oct 2026 (FYP-1, to mid-term) & Base architecture with implementation and performance benchmark; crop-type model evaluation and comparison.\\
Iteration 2 & mid-Oct -- Dec 2026 (FYP-1, to final) & Integrate the selected crop-type pipeline into the platform; CDL confidence rule; cross-area test; multi-image upload.\\
Iteration 3 & Feb -- Mar 2027 (FYP-2) & Field-boundary model and area validation; yield data and first yield model; more AOIs.\\
Iteration 4 & Apr -- May 2027 (FYP-2) & End-to-end evaluation, hardening, deployment and final documentation.\\
\bottomrule
\end{tabularx}
\end{table}

\section{Iteration 1 Plan}\label{sec:it1plan}
\textbf{Objectives.} (1) Replace the prototype with a base architecture that stays fast as images and
models grow, implemented and benchmarked on real data. (2) Find a crop-type model that is accurate
under a fair test and light enough for the platform.

\textbf{Selected features.} Tile-based Live Map with layer panel, opacity, legends, crop filter,
field search and point inspection; dashboard statistics from the database; upload and background
processing with progress; model registry with the current cropland and crop-type models;
system-status page. On the model side: monthly composites for AOI1 and AOI2, eleven-model
comparison, improvement ladder, error analysis and comparison with published papers.

\textbf{Deliverables.} Platform repository with Docker Compose stack and documentation; benchmark
results; crop-type results report; this mid-report.

\textbf{Evidence.} Benchmark tables, test results, screenshots (Section~\ref{sec:demo}) and
repository history (Appendix~\ref{app:repo}).

\section{Iteration 2 Plan}\label{sec:it2plan}
\textbf{Objectives.} Bring the selected crop-type pipeline into the platform and strengthen the
evaluation.

\textbf{Selected features.}
\begin{itemize}
  \item Package the LightGBM pipeline (monthly composites, neighbourhood inputs, 3$\times$3 vote) as a
        registry model and replace the baseline crop-type model in the platform.
  \item Monthly-composite ingestion in the processing pipeline.
  \item Fix one CDL confidence rule and report it next to the all-pixel score; compute the AOI2
        interval and interior table.
  \item Cross-area test (train on one AOI, test on the other).
  \item Multi-image upload and an ``all areas'' map view.
  \item Measure end-to-end serving cost of the improved pipeline on CPU and repeat the workflow
        resource measurement (Table~\ref{tab:resources}).
  \item Usability study: task-based sessions with a questionnaire score and task times (P3).
  \item Multi-layer test: eight or more layers visible at once (NFR-P3).
\end{itemize}
\textbf{Deliverables.} Updated platform release, evaluation addendum and FYP-1 final report.
\textbf{Evidence.} Test results, benchmark re-run and model cards with stored metrics.

\section{Work Breakdown and Responsibilities}
\begin{table}[htbp]
\centering\small
\caption{Work breakdown and ownership}
\label{tab:wbs}
\begin{tabularx}{\textwidth}{L{3.4cm} Y L{3.5cm}}
\toprule
\textbf{Work package} & \textbf{Tasks} & \textbf{Owner}\\
\midrule
WP1 Architecture and platform & Architecture design; performance analysis of the prototype; tile serving; storage and database design; processing pipeline; model registry; web portal; deployment configuration; benchmark. & Huzaifa Bin Khalid\\
WP2 Crop-type model evaluation & Monthly composite preparation; spatial block protocol; model comparison and cost measurement; improvement ladder; error analysis. & Saneed Ullah, Taha Bashir Awan\\
WP3 Comparison with literature & Reading and checking reference papers; like-for-like comparison and limits. & Saneed Ullah, Taha Bashir Awan\\
WP4 Integration (Iteration 2) & Packaging the selected model as a registry plug-in; end-to-end measurement. & Shared\\
WP5 Documentation and evaluation & Reports, demonstrations, final evaluation. & Shared\\
\bottomrule
\end{tabularx}
\end{table}

\section{Progress Against Plan}
\begin{table}[htbp]
\centering\small
\caption{Planned phases (proposal slides) against actual progress}
\label{tab:progress}
\begin{tabularx}{\textwidth}{L{3.3cm} L{4.5cm} Y}
\toprule
\textbf{Planned (slides)} & \textbf{Planned content} & \textbf{Actual at mid-term}\\
\midrule
Phase 1 (FYP-1) Research and data understanding & Literature review; CDL/AOI analysis; satellite data acquisition & Done. Literature extended with eight crop-type papers; monthly 2022 composites exported for AOI1 and AOI2.\\
Phase 2 (FYP-1) Core crop mapping pipeline & Preprocessing; training dataset; baseline model; initial GIS boundary extraction & Done for two AOIs. Eleven models compared; improved pipeline selected. Field regions derived from crop-type output as interim boundaries.\\
Phase 3 (FYP-2) Yield and platform integration & Yield module; Web GIS; system integration & Platform architecture and integration brought forward and done in FYP-1. Yield remains in FYP-2.\\
Phase 4 (FYP-2) Evaluation and finalisation & Hold-out tests; failure analysis; performance evaluation; deployment & Partly brought forward: spatial hold-out testing, error analysis and performance benchmark already done; final evaluation and deployment remain.\\
\bottomrule
\end{tabularx}
\end{table}

\section{Delays, Deviations and Recovery Plan}
\begin{itemize}
  \item \textbf{Platform work brought forward (deviation).} The slides placed Web GIS and integration in
        FYP-2. Profiling showed the prototype would not scale, and the panel named performance,
        light weight, usability and integration as the main concerns (P3--P6), so the architecture
        was redesigned and implemented in Iteration 1. This reduces
        integration risk in FYP-2.
  \item \textbf{Technology change (deviation).} The Django and Leaflet prototype that sent GeoJSON was
        replaced by FastAPI, React with MapLibre, TiTiler, PMTiles, PostGIS, Celery and MinIO
        (Section~\ref{sec:evolution}).
  \item \textbf{Model recommendation reversed (deviation).} Early research notes recommended the
        Presto foundation model; the experiments did not support it for the U.S./CDL case, so a
        LightGBM pipeline is proposed.
  \item \textbf{Platform still runs the baseline crop-type model (delay).} The selected pipeline needs
        monthly composites at ingest; integration is the first item of Iteration 2.
  \item \textbf{Boundary model not started (delay).} Field regions derived from crop-type output act as
        the interim boundary layer; a boundary model is planned for Iteration 3.
  \item \textbf{Two of eight AOIs processed (delay).} Monthly composites exist for AOI1 and AOI2; the
        remaining AOIs will be exported for the cross-area and generalisation tests.
  \item \textbf{AI-assisted implementation (deviation from the stated GenAI boundary).} The proposal
        limited generative AI to documentation; during Iteration 1 an AI coding assistant was also
        used for parts of the platform implementation, under the team's design and review
        (Section~\ref{sec:tools-assisted}).
\end{itemize}
\textbf{Recovery.} Iteration 2 starts with model integration, then the CDL confidence rule and the
cross-area test; AOI exports run in the background on Google Earth Engine. The boundary and yield
work in FYP-2 benefits from the registry and pipeline already in place, so new models only need a
task and a manifest.
